\section{All game instructions}
\label{app:instruct}

\noindent\textbf{Basic 11-20 Game}
\begin{quote}
\emph{You and another player are playing a game in which each player requests an amount of money. The amount must be (an integer) between 11 and 20 shekels.  Each player will receive the amount he requests. A player will receive an additional amount of 20 shekels if he asks for exactly one shekel less than the other player. What amount of money would you request?}
\end{quote}

\noindent\textbf{Cycle 11-20 Game}
\begin{quote}
\emph{You and another player are playing a game in which each player requests an amount of money.
The amount must be (an integer) between 11 and 20 shekels. 
Each player will receive the amount of money he requests.
A player will receive an additional amount of 20 shekels if: (i)  he asks for exactly one shekel less than the other player or  (ii) he asks for 20 shekels and the other player asks for 11 shekels.
What amount of money would you request?}
\end{quote}

\noindent\textbf{Costless 11-20 Game}
\begin{quote}
\emph{You and another player are playing a game in which each player chooses an integer in the range 11-20.
A player who chooses 20 will receive 20 shekels (regardless of the other player's choice). 
A player who chooses any other number in this range will receive three shekels less than in the case where he chooses 20. 
However, he will receive an additional amount of 20 shekels if he chooses a number that is one less than that chosen by the other player. 
Which number would you choose?}
\end{quote}

\noindent\textbf{Basic 1-10 Game}
\begin{quote}
\emph{You are going to play a game where you must select a whole number between 1 and 10.
You will receive a number of points equivalent to that number. For example, if you select 3, you will get 3 points. 
If you select 7, you will get 7 points, etc.
After you tell us your number, we will randomly pair you with another Prolific worker who is also playing this game. 
They will also have chosen a number between 1 and 10.
If either of you select a number exactly one less than the other player's number, than the player with the lower number will receive an additional 10 points.
Please choose a number between 1 and 10. }
\end{quote}

\noindent\textbf{Cycle 1-10 Game}
\begin{quote}
\emph{You are going to play a game where you must select a whole number between 1 and 10.
You will receive a number of points equivalent to that number. 
For example, if you select 3, you will get 3 points. 
If you select 7, you will get 7 points, etc.
After you tell us your number, we will randomly pair you with another Prolific worker who is also playing this game. 
They will also have chosen a number between 1 and 10.
There are 2 ways to win an additional 10 points based on both yours and the other player's choice:
1. If either of you select a number exactly one less than the other player's number, then the player with the lower number will receive an additional 10 points.
2. If either of you select 10 and the other selects 1, then the player who chose 10 will receive an additional 10 points.
Please choose a number between 1 and 10.}
\end{quote}

\noindent\textbf{Costless 1-10 Game}
\begin{quote}
\emph{You are going to play a game where you must select a whole number between 1 and 10.
You will receive 10 points if you select the number 10 and you will receive 7 points for selecting any other number. 
After you tell us your number, we will randomly pair you with another Prolific worker who is also playing this game. 
They will also have chosen a number between 1 and 10.
If either of you select a number exactly one less than the other player's number, than the player with the lower number will receive an additional 10 points.
Please choose a number between 1 and 10.}
\end{quote}

\noindent\textbf{1-7 Game}
\begin{quote}
\emph{You are going to play a game where you must select a whole number between 1 and 7.
You will receive a number of points equivalent to that number. For example, if you select 3, you will get 3 points. 
If you select 6, you will get 6 points, etc.
After you tell us your number, we will randomly pair you with another Prolific worker who is also playing this game. 
They will also have chosen a number between 1 and 7.
If either of you select a number exactly one less than the other player's number, than the player with the lower number will receive an additional 10 points.
Please choose a number between 1 and 7.}
\end{quote}

\begin{figure}[h!]
\begin{center}
\begin{minipage}{1 \linewidth}
\caption{Screenshot of the three-player game instructions} 
\centering
\vspace*{-0.15in}
\label{fig:novel_allocation_pic}
\includegraphics[width=\textwidth]{images/novel_allocation.png}
\end{minipage}
\end{center}
\begin{footnotesize}
\begin{singlespace}
\vspace*{-0.15in}
\emph{Notes:} This figure shows the instructions for the novel three-player allocation game presented to participants.
\end{singlespace}
\end{footnotesize}
\end{figure}

\begin{figure}[h!]
\begin{center}
\begin{minipage}{1 \linewidth}
\caption{Bonus opportunity for the games in Section~\ref{sec:ext-valid}} 
\centering
\vspace*{-0.15in}
\label{fig:bonus_bounds_pic}
\includegraphics[width=.75\textwidth]{images/bonus_bounds.png}
\end{minipage}
\end{center}
\begin{footnotesize}
\begin{singlespace}
\vspace*{-0.15in}
\emph{Notes:} This shows the instructions for the bonus opportunity presented to participants for the novel sample of 1,500 games.
\end{singlespace}
\end{footnotesize}
\end{figure}

\newpage \clearpage

\begin{figure}[h!]
\begin{center}
\begin{minipage}{1 \linewidth}
\caption{Choosing a number for the assigned game in Section~\ref{sec:ext-valid}} 
\centering
\vspace*{-0.15in}
\label{fig:choice_bounds_pic}
\includegraphics[width=0.75\textwidth]{images/choice_bounds.png}
\end{minipage}
\end{center}
\begin{footnotesize}
\begin{singlespace}
\vspace*{-0.15in}
\emph{Notes:} This figure shows an example screenshot of participants selecting their number for their assigned game from the set of 1,500 games.
\end{singlespace}
\end{footnotesize}
\end{figure}

\renewcommand{\thefigure}{C\arabic{figure}} 
\setcounter{figure}{0}  

\renewcommand{\thetable}{C\arabic{table}} 
\setcounter{table}{0}  
